% !TeX root = ../thuthesis-example.tex
\chapter{安装及使用}
\section{系统运行环境}
系统支持在电脑网页和移动设备上运行。网页端可通过主流浏览器（如 Chrome、Firefox）直接访问，无需安装；移动端提供适配 Android 和 iOS 的安装包，用户可通过官网下载并安装。

平台需全程联网运行，语音识别与智能分析均在云端完成。为了保证音频上传和分析的速度，建议使用稳定的 Wi-Fi 或 5G 网络。

设备需配备可正常工作的麦克风。第一次使用时，需允许 App 或浏览器访问录音权限，以便系统采集语音数据。为了获得更高的识别准确度，推荐使用带有清晰录音功能的设备。

注册方式采用手机号验证，用户输入手机号并获取验证码即可完成注册登录。对于想先体验的用户，平台提供游客模式，无需注册即可直接进入练习。不过游客模式下的数据不会被长期保存，退出后将被清空。

在移动端安装时，Android 用户需开启“允许未知来源”才能安装 APK 文件；iOS 用户则可通过app store下载。网页端无需安装，直接在浏览器输入平台网址即可进入主界面。

\section{系统主要功能}
本系统面向语言发音学习，结合了语音识别、知识图谱和大模型智能分析等技术，为用户提供发音练习、错误原因分析、智能学习规划以及知识图谱浏览等功能。以下将分别介绍各主要功能模块的使用方法和特点。

\subsection{发音练习模块}
发音练习是本系统的核心功能模块。如\ref{fig:发音界面}所示，当用户从主界面点击“学习”进入发音练习界面后，系统将显示本次练习的朗读词语。用户可以先阅读内容并点击“播放”按钮，系统会播放词语标准读法。用户点击开始录制后，系统会提示“录制”，用户开始朗读给出的内容。用户每次练习可以读三次，朗读结束后点击“停止”或等待系统自动结束录音。若对自己的朗读不满意，可再次点击录音按钮重读内容。系统会保留最近的三次录音，并选择最优进行供分析。
\begin{figure}[H]
    \centering
    \includegraphics[scale=0.25]{发音界面.png}
    \caption{发音界面}
    \label{fig:发音界面}
\end{figure}
 用户提交录音后，系统将调用云端服务对语音进行声学和语音学分析。科大讯飞“星火”认知大模型提供了强大的跨领域语言理解与推理能力 。借助其对方言的深度理解，方言声学知识图谱和现代汉语方言音档，系统后台的智能体会对用户的发音进行全面分析，包括语音的准确度、清晰度、语调语速等方面。系统采用了多智能体协同架构，例如由“发音教练Agent”分析用户发音准确性，“语音韵律Agent”评估语调和流利度等，各智能体分工合作对语音从不同维度实时解析 。分析过程在数秒内完成，并会自动生成语图对比结果与评测反馈，如\ref{fig:语图界面}和\ref{fig:语图分析}所示，清晰标出发音偏差和改进方向。这些可视化结果与文字建议会即时呈现在用户端，方便学习者直观了解自己的表现，并作为后续训练和智能规划的依据
\begin{figure}[H]
    \centering
    \includegraphics[scale=0.25]{语图界面.png}
    \caption{语图界面}
    \label{fig:语图界面}
\end{figure}
\begin{figure}[H]
    \centering
    \includegraphics[scale=0.25]{语图分析.png}
    \caption{语图分析界面}
    \label{fig:语图分析}
\end{figure}

\subsection{错因分析与智能规划}
每次练习结束后，用户进入“练习基地”模块查看自己的发音分析报告和系统生成的学习规划。练习基地相当于个人学习中心，记录并展示用户历次练习的结果数据和改进建议，其主要功能包括错因分析和智能规划两部分.
\paragraph{错因分析}
如图\ref{fig:错因分析}所示，系统会利用内置的方言声学知识图谱，语言模型和现代汉语音档，对错误原因进行智能推断。如某用户在普通话练习中某些时长节奏或者声调出现问题，系统会指出这一情况并解释可能的成因。通过这种分析，用户可以清楚了解自己发音问题的症结所在。系统拥有强大的识别能力，能够突破方言口音的干扰精准识别细微发音偏误 。报告将错误音素的位置和程度直观地展示给用户 。借助这些多维度的评估结果，用户获得了类似专业老师批改的详尽反馈 。
\paragraph{智能规划}
如图\ref{fig:智能规划}所示，在完成错误分析的基础上，系统的AI智能体会根据用户当前的发音水平和错误模式，自动制定个性化的学习规划建议 。智能规划以文本形式呈现在练习基地页面，内容包括针对用户弱项的强化练习建议和下阶段学习目标。例如，若用户多次练习中都出现某个音素发音不准，智能规划会建议重点练习该音素，并列出相应的示例字词或绕口令供下次练习使用。又如针对声调错误，系统会规划一系列同声调对比练习或节奏训练来帮助纠正。系统会根据智能规划为用户生成下次练习的题目，做到有的放矢。例如，上例中的用户在下次进入练习时，将优先收到包含第三声的练习词语，以强化该弱项。如此形成“练习-反馈-再练习”的闭环，使得学习路径根据用户表现自适应调整，真正实现个性化的发音教学。
 \begin{figure}[H]
    \centering
    \includegraphics[scale=0.25]{错因分析.png}
    \caption{错因分析界面}
    \label{fig:错因分析}
\end{figure}
\begin{figure}[H]
    \centering
    \includegraphics[scale=0.25]{智能规划.png}
    \caption{智能规划界面}
    \label{fig:智能规划}
\end{figure}
\paragraph{知识图谱浏览}
本系统独特地提供了知识图谱浏览功能，用户可以通过该模块直观地查看支撑系统的方言声学知识图谱内容。知识图谱以图形可视化界面呈现，将语音领域的相关知识以节点和连线的形式展示出来，帮助用户更好地理解发音要点和方言差异。

进入知识图谱浏览界面后，用户可以看到若干节点，例如音素/音标、音系特征、方言、示例词汇等核心实体节点。节点之间由连线连接，表示它们之间的关系。例如，某个音素节点连接到不同方言节点，表示该音素在这些方言中的对应发音；音系特征节点连接到具体音素，表示该音素具有的语音特征。用户可以拖动、放大缩小图谱视图来浏览不同区域的内容。当点击某一节点时，界面会显示该节点的详细信息。例如，点击一个音素节点，可能显示该音素的国际音标符号、发音部位和方法说明，还有该音素在普通话及若干地方方言中的发音差异说明。若节点有关联的多媒体信息，用户还可能播放对应的音频样例或查看口型图片等。比如点击某方言节点，界面可能提供该方言代表性语音片段的播放，让用户直观感受该方言的语音特点。

知识图谱浏览功能不仅是一个学习工具，也体现了系统的技术支撑体系。该知识图谱是基于海量语音和文本数据构建而成，涵盖了普通话及各主要方言的声学知识。通过对多源异构的语料进行采集和强制对齐，以及自动化的知识抽取与融合，系统构建了完善的声学知识库 。图谱经过质量评估，保证了其准确性和实用性，并在系统核心算法中发挥重要作用。对于用户而言，通过浏览知识图谱，可以深化对语言发音规律的理解。例如，用户可以查询某个自己总发不准的音素在知识图谱中的位置，从中发现其发音要领或者与其他音素的区别；也可以通过图谱了解某地方言与普通话的发音对应关系，满足好奇心的同时辅助纠正方音。有了知识图谱的直观呈现，抽象的语音概念变得更易理解，学习者可以将零散的发音知识串联成系统的框架，加深记忆和应用。
\begin{figure}[H]
    \centering
    \includegraphics[scale=0.25]{知识图谱.png}
    \caption{知识图谱界面}
    \label{fig:知识图谱}
\end{figure}
